\documentclass[runningheads]{llncs}
\usepackage[T1]{fontenc}
\usepackage{graphicx}
\usepackage{float}
\usepackage{booktabs}
\usepackage{tabularx}
\usepackage{color}
\usepackage[colorlinks=true, linkcolor=black, citecolor=black, urlcolor=blue]{hyperref}
\renewcommand\UrlFont{\color{blue}\rmfamily}
\urlstyle{rm}
\usepackage{soul}
\begin{document}
\title{MSc Project 1: Time-Series-Enhanced Predictive Process Monitoring}
\subtitle{Phase Review 4}
\titlerunning{Time-Series-Enhanced PPM: Phase Review 4}
\author{Micha\l{} Durkalec \and
    Konstantinos Sakkas \and
    Roman Le\v{z}ovi\v{c}
}
\authorrunning{M. Durkalec et al.}
\institute{Supervised Process Intelligence (SPI) Lab,\\
    RWTH Aachen University, Aachen, Germany\\
    \email{office@pads.rwth-aachen.de}}
\maketitle
\begin{abstract}
    This document reviews the Testing and Assessment phase of the project. Building on the
    infrastructure completed in Sprint~3, this phase ran the full experimental comparison
    between log-only and time-series-enhanced models for both classification and regression
    tasks on the complete BPI Challenge 2017 dataset, implemented explainability and
    diagnostic tooling, and produced the empirical results that answer the project's
    central research question. The phase concludes with all experimental findings
    documented and the final report as the remaining deliverable.
    \keywords{Predictive Process Monitoring \and Model Evaluation \and Feature Importance
        \and Time-Series Features \and Overfitting Analysis}
\end{abstract}

\section{Introduction \& Objectives}

\subsection{Phase Goals}
The Testing and Assessment phase aimed to run the full experimental suite on the complete
BPI Challenge 2017 dataset~\cite{bpic2017}, moving beyond the development subset used
during earlier sprints. We compared log-only and time-series-enhanced models for both
outcome classification and remaining-time regression under identical conditions,
producing the head-to-head comparison deferred from Sprint~3. We also built diagnostic
and explainability tooling: permutation feature importance, SHAP summary plots,
overfitting analysis, confusion matrices, and predicted-versus-actual scatter plots to
interpret the models and assess their reliability.

\subsection{Team Organization}
Work continued to be tracked through GitHub Issues. Michal drove most of the
experimental pipeline: model comparisons for both classification and regression,
confusion matrices, overfitting metrics, and feature importance data.
Konstantinos built the SHAP explainability module, added binary outcome
classification, refactored the evaluation metrics for the final report, and wrote the
architecture documentation. Roman handled class-balanced training with
cross-validation search, implemented the baseline models, and fixed the prefix generation
logic for the classification task. All team members contributed to refining the empirical
findings.

\subsection{Role Allocation}
Role assignments from Sprint~3 carried forward into this phase, with all team members
contributing across tracks as needed.

\begin{table}[H]
    \caption{Role assignments for the Testing and Assessment phase.}\label{tab:roles}
    \centering
    \begin{tabular}{@{}ll@{}}
        \toprule
        \textbf{Role}                      & \textbf{Assigned To} \\
        \midrule
        Project Coordinator / Scrum Master & Roman                \\
        Technical Lead / Architecture      & Konstantinos         \\
        Data \& Preprocessing Lead         & Roman                \\
        Modelling \& Experimentation Lead  & Michal               \\
        Evaluation \& Visualization Lead   & Konstantinos         \\
        Quality \& Tooling (Tests, CI)     & Roman                \\
        Documentation \& Reporting         & Michal, Roman, Konstantinos (shared) \\
        \bottomrule
    \end{tabular}
\end{table}

\section{Technical Implementation}

This phase did not introduce new architecture. We ran the existing
infrastructure from Sprints~1--3 at full scale on the complete BPI Challenge
2017 dataset. \autoref{tab:steps} maps each work item
to its issues and pull requests.

\subsection{Full-Scale Experimental Runs}
We ran all experiments through the parametric CLI entry point from Sprint~2.
For classification we compared four models: Dummy, Logistic Regression, Random Forest,
and HistGradientBoosting using the full 58-feature log+time-series set. For regression
Ridge replaced Logistic Regression in the same four-model design. All
configurations shared the 0.8-quantile temporal train/test split by case start time, prefix
lengths from~3 to~100, and identical random seeds. Hyperparameter search used
\texttt{RandomizedSearchCV} with 3-fold cross-validation and 10 iterations. HGB came out
best for classification; Random Forest for regression.

\subsection{Time-Series Impact Analysis}
For the central comparison we ran the best model from each task under both configurations
with all other settings held constant. Hourly-resampled time-series features with
8-hour rolling windows did not improve classification and degraded regression
performance compared to the log-only baseline. Unbounded prefixes and capped prefixes
produced nearly identical results.

\subsection{Diagnostic and Explainability Tooling}
We ran several diagnostic analyses to interpret the models. Permutation feature
importance ranked activity-count features at the top for both tasks. SHAP summary plots
visualised how individual feature values push predictions toward specific outcomes.
Overfitting looked moderate: train-test gaps stayed within expected bounds for tree-based
ensembles. We also produced confusion matrices and predicted-versus-actual scatter plots
for every model-task combination.

\begin{table}[H]
    \centering
    \caption{Testing and Assessment implementation steps.}
    \label{tab:steps}
    \begin{tabularx}{\textwidth}{lXX}
        \toprule
        \textbf{Step} & \textbf{Description}                                   & \textbf{Issues / PRs} \\
        \midrule
        1             & Run experiments for both tasks                          & \#120                 \\
        2             & Regression model comparison                             & \#119                 \\
        3             & Classification model comparison with HGB                & (\#112--\#118)        \\
        4             & Confusion matrices                                      & --                    \\
        5             & Overfitting metrics                                     & --                    \\
        6             & Feature importance to CSV                               & --                    \\
        7             & Predicted-vs-actual scatter plots                       & --                    \\
        8             & HGB as best classification model                        & --                    \\
        9             & Full-prefix-length configuration                        & --                    \\
        10            & Test suite maintenance and fixes                        & --                    \\
        \bottomrule
    \end{tabularx}
\end{table}

\section{Review \& Retrospective}

\subsection{Successes}
All planned phase objectives were completed. The full experimental suite ran end-to-end
on the complete BPI Challenge 2017 dataset, producing the head-to-head comparison
deferred from Sprint~3. The diagnostic tooling ran without issues and its results fed
into the interpretation. \autoref{tab:issues} lists the key issues and pull
requests closed during the phase.

\begin{table}[H]
    \centering
    \caption{Testing and Assessment issues and pull requests.}
    \label{tab:issues}
    \begin{tabularx}{0.85\textwidth}{lX}
        \toprule
        \textbf{Issue / PR} & \textbf{Description}                                            \\
        \midrule
        \#120               & Run experiments for both tasks                                  \\
        \#119               & Compare regression models                                       \\
        \#112--\#118        & Classification experiments with HGB                              \\
        --                  & Comput confusion matrices                                       \\
        --                  & Evaluate overfitting metrics                                    \\
        --                  & Save feature importance to CSV                                  \\
        --                  & Predicted vs actual scatter plots                               \\
        --                  & HGB as best model for classification                            \\
        --                  & Add config for full prefix len                                  \\
        --                  & Fix tests (remove integration tests)                            \\
        \bottomrule
    \end{tabularx}
\end{table}

\subsection{Difficulties}
The main challenge was the negative result itself: the time-series features added no
value for classification and degraded regression. Explaining this outcome required
deeper diagnostic analysis than initially scoped. The key insight was concept drift:
the time-series features are global process-state signals whose distributions shift
between train and test on the chronological split, so they fail to generalise.

\subsection{Task Tracking}
We managed all work through GitHub Issues with linked pull requests. Smaller changes
were merged directly to dev, while larger features followed the standard review
process. The requirements tracking system kept detailed sub-requirements as GitHub
issues linked to parent requirements instead of editing the requirements document,
which worked well for scope changes like adding HGB as a model family.

\section{Moving Forward}

\subsection{Future Planning}
With the Testing and Assessment phase complete, the remaining deliverable is the final
report. We will produce the document covering data preprocessing, feature
engineering, model specification, experimental protocol, and results. It will include a
critical evaluation of the time-series features, limitations, and a project retrospective.
All remaining open issues will be closed and the repository documentation finalised
alongside.

\subsection{Lessons Learned}
The main lesson from this phase is that building extraction machinery is not the same as
validating that the features carry useful signal. We put more effort into constructing the
pipeline than into checking whether the extracted signals tracked real workload variation,
and the null result reflects that imbalance. Claims about time-series features also need
to be scoped to the specific extraction strategy rather than generalised. These are
observations we carry forward into the final report.

\begin{thebibliography}{8}
    \bibitem{bpic2017}
    van Dongen, B.F.: BPI Challenge 2017. Eindhoven University of Technology (2017). \url{https://research.tue.nl/en/datasets/bpi-challenge-2017}
\end{thebibliography}
\end{document}
